% What is on the bench, and what a real safety chain has that this one does not.
\begin{tikzpicture}[font=\sffamily\scriptsize, x=1cm, y=1cm]

\node[chlabel, anchor=west] at (-6.8,3.4) {what is on the bench, and what an integrity claim would need beside it};

% ================================================================ the real chain
\node[layerlabel] at (-6.8,2.9) {single channel, and all of it real};
\node[hw, text width=26mm, minimum height=9mm] (sen) at (-5.2,2.0)
  {the ranging shield\\8 by 8 zones};
\node[fw, text width=26mm, minimum height=9mm] (mcu) at (-1.8,2.0)
  {the node: machine and latch};
\node[hw, text width=26mm, minimum height=9mm] (pin) at (1.6,2.0)
  {the break input, then one output pin};
\node[ledsym, fill=red!30, minimum size=5mm] (led) at (3.8,2.0) {};
\draw[flow] (sen) -- (mcu);
\draw[faultedge] (mcu) -- (pin);
\draw[faultedge] (pin) -- (led);
\node[chlabel, anchor=west] at (4.2,2.0) {the indicator, which is all the reaction there is};

% ================================================================ what is absent
\node[layerlabel] at (-6.8,0.9) {and what is not here, drawn rather than implied};
\node[absentblk, text width=26mm, minimum height=9mm] (ch2) at (-5.2,0.0)
  {a second channel, and a comparison};
\node[absentblk, text width=26mm, minimum height=9mm] (wdg) at (-1.8,0.0)
  {an external watchdog};
\node[absentblk, text width=26mm, minimum height=9mm] (mot) at (1.6,0.0)
  {a motor and a brake};
\node[absentblk, text width=26mm, minimum height=9mm] (ass) at (5.0,0.0)
  {an assessment, and a certificate};

\node[chlabel, anchor=north, text width=30mm, align=center] at (-5.2,-0.6)
  {without it, the architecture stops at the lower levels by construction};
\node[chlabel, anchor=north, text width=30mm, align=center] at (-1.8,-0.6)
  {both on-die watchdogs share the chip's fate};
\node[chlabel, anchor=north, text width=30mm, align=center] at (1.6,-0.6)
  {so the second safe state is designed and never felt};
\node[chlabel, anchor=north, text width=30mm, align=center] at (5.0,-0.6)
  {and no integrity level is claimed anywhere};

% ================================================================ the arithmetic
\node[regcell, fill=hwcol, minimum width=114mm, text width=112mm, inner sep=3pt,
      anchor=north west, align=left] at (-6.8,-2.2)
  {\textbf{The architecture arithmetic, in one line each, from a free vendor
   guide:} dual channel with hardware fault tolerance one reaches the higher
   integrity and performance levels; \textbf{single channel with fault tolerance
   zero does not}, whatever the quality of the firmware. That is architecture
   rather than effort, and it is why this chapter's claim is worded the way it
   is.};

\node[umlnote, text width=114mm, anchor=north west] at (-6.8,-3.6)
  {The same guide carries the distinction firmware engineers most often miss:
   \textbf{random hardware capability and systematic capability are different
   things}, rated separately on the same part. Systematic capability is about
   the development process; random hardware capability is about diagnostic
   coverage and failure rates. Most parts are rated differently on each, and a
   single number quoted without saying which one it is means very little.};
\end{tikzpicture}
